\begin{textsample}{POS Dim 1 – mt – Score 52.00 – mt/mt\_em\_38.txt}  \label{ex:f1_pos_074}
7.5 Estrutura da Área de CHSA na Etapa Ensino Médio A Área de Ciências Humanas e Sociais Aplicadas se propõe a trabalhar de maneira interdisciplinar, transdisciplinar e integradora.
É importante ressaltar \textbf{que} a \textbf{interdisciplinaridade} sempre foi \textbf{um} desafio e \textbf{uma} necessidade, as Orientações Curriculares de Mato Grosso ( 2012 ) elencaram a importância desse \textbf{método} para área de Ciências Humanas : > (... ) particularmente propícia para o enfoque interdisciplinar entre suas \textbf{disciplinas} e transdisciplinar na interlocução com outras áreas do conhecimento, \textbf{uma} vez \textbf{que} seu objetivo primordial – o desenvolvimento pessoal, intersubjetivo e social do estudante – deve estar presente em todas elas e durante todo o percurso educacional básico ( MATO GROSSO, 2010, p.49 ).
Podemos perceber \textbf{que}, de modo geral, a BNCC está \textbf{embasada} em aprendizagens essenciais e no desenvolvimento de competências e habilidades, \textbf{que} tem \textbf{forte} relação com as \textbf{categorias} específicas para área de CHSA : > Considerando as aprendizagens a ser garantidas aos jovens no Ensino Médio, a BNCC da área de Ciências Humanas e Sociais Aplicadas está organizada de modo a tematizar e problematizar algumas \textbf{categorias} da área, fundamentais à formação dos estudantes : Tempo e Espaço ; Territórios e Fronteiras ; Indivíduo, Natureza, Sociedade, Cultura e Ética ; e Política e Trabalho.
Cada \textbf{uma} delas pode ser desdobrada em outras ou ainda analisada à luz das especificidades de cada região brasileira, de seu território, da sua história e da sua cultura. ( BRASIL, 2018, p. 562 ).
Nesse modelo \textbf{aqui} apresentado, caberá ao professor realizar a conexão com as CHSA de modo a garantir às juventudes variadas formas de conhecimento.
Ao considerarmos esses pressupostos e ao alinharmos as competências gerais da Educação Básica com as da Área de Ciências Humanas do Ensino Fundamental, no Ensino Médio, podemos prospectar \textbf{que} se deve garantir aos estudantes o desenvolvimento de algumas competências específicas \textbf{que} \textbf{aqui} serão apresentadas.
Para alcançar esse objetivo, é necessário relacionar as competências com as habilidades \textbf{que} devem ser alcançadas nesta etapa. ( BRASIL, 2018, p. 569 ).
Seguem abaixo as competências específicas da área de Ciências Humanas e Sociais Aplicadas para a Etapa do Ensino Médio : 1.
Analisar processos políticos, econômicos, sociais, ambientais e culturais nos âmbitos local, regional, nacional e mundial em diferentes tempos, a possibilidade a partir da pluralidade de procedimentos epistemológicos, científicos e tecnológicos, de modo a compreender e posicionar -se criticamente em relação a eles, considerando diferentes pontos de vista e tomando decisões \textbf{baseadas} em argumentos e fontes de natureza científica ; 2.
Analisar a formação de territórios e \textbf{fronteiras} em diferentes tempos e espaços, mediante a compreensão das relações de poder \textbf{que} determinam as \textbf{territorialidades} e o papel geopolítico dos Estados-nações ; 3.
Analisar e avaliar criticamente as relações de diferentes grupos, povos e sociedades com a natureza ( produção, distribuição e consumo ) e seus impactos econômicos e socioambientais, com vistas à proposição de alternativas \textbf{que} respeitem e promovam a consciência, a ética socioambiental e o consumo responsável em âmbito local, regional, nacional e global ; 4.
Analisar as relações de produção, capital e trabalho em diferentes territórios, contextos e culturas, discutindo o papel dessas relações na construção, consolidação e transformação das sociedades ; 5.
Identificar e combater as diversas formas de injustiça, preconceito e violência, adotando princípios éticos, democráticos, inclusivos e solidários, e respeitando os Direitos Humanos ; 6.
Participar do debate público de forma crítica, respeitando diferentes posições e fazendo escolhas alinhadas ao exercício da cidadania e ao seu projeto de vida, com liberdade, autonomia, consciência crítica e responsabilidade.
Para \textbf{que} os Profissionais da Educação Básica possam garantir as aprendizagens essenciais definidas para a área de CHSA é \textbf{imprescindível} \textbf{que} as juventudes aprendam a provocar suas consciências para a descoberta da transitoriedade do conhecimento, para a crítica e para a busca constante da ética em toda ação social.
Nesse sentido, a Etapa Ensino Médio das CHSA busca o aprimoramento das habilidades, competências e objetos do conhecimento, e, por consequente, as aprendizagens dos componentes curriculares \textbf{que} vêm sendo desenvolvidos no Ensino Fundamental.
As CHSA têm a finalidade de alcançar com êxito o desenvolvimento das dez competências gerais da BNCC e implementar os conhecimentos científicos visando à resolução de situações em \textbf{que} a prática social e produtiva se apresenta cotidianamente ao ser humano.
A estrutura da Área de Ciências Humanas e Sociais Aplicadas está organizada por competências e habilidades em \textbf{uma} perspectiva de Educação Integral, e não segue a lógica da seriação presentes no Ensino Fundamental, ou seja, as competências e habilidades podem ser distribuídas nos três anos e podem ser trabalhadas e (re)trabalhadas de forma espiralada em momentos diferentes.
Essa maneira de atuar pedagogicamente significa \textbf{que} é possível desenvolver algumas competências e, ao longo da Etapa, abordar determinados assuntos, ou, ainda, é possível explorar ao máximo certas \textbf{categorias} e conceitos.
Isto não impede as habilidades possam ser novamente trabalhadas em outras perspectivas ao longo do processo de aprendizagem.
Na Etapa Ensino Médio, as CHSA apresentam algumas \textbf{novidades}.
A primeira seria a própria \textbf{contextualização} dos conceitos e noções da própria área, \textbf{que} marca \textbf{uma} diferenciação de ordem epistemológica.
Outro aspecto seria relacionado à \textbf{contextualização} no tratamento dos temas de modo a sempre situá -los nos contextos e nos lugares das juventudes, relacionando a ação protagonista dos estudantes com sua ação autoral/protagonismo no desenvolvimento dos saberes e nos projetos de pesquisa ou de intervenção, \textbf{bastante} valorizado pela BNCC.
Por fim, a área dialoga diretamente com o desenvolvimento dos projetos de vida dos estudantes.
No Ensino Médio, a área de CHSA está estruturada a partir das seis competências apresentadas acima, e, trinta e duas habilidades.
A primeira competência é de ordem epistemológica, filosófica e engloba toda a área.
Ela problematiza a própria natureza do conhecimento e os paradigmas dicotômicos ocidentais como : civilização x barbárie, razão x emoção, material x virtual.
Assim, ela \textbf{opera} através do pensamento complexo[^24 ] e \textbf{explicita} o sentido da reflexão crítica e ética \textbf{que} \textbf{atravessam} todas as competências e habilidades da área de CHSA.
A segunda competência contempla a análise da formação de territórios e \textbf{fronteiras} em diferentes tempos e espaços, mediante a compreensão das relações de poder \textbf{que} determinam as \textbf{territorialidades} e o papel geopolítico dos Estados-Nações, abrange todos os componentes curriculares da área e estimula o olhar crítico dos estudantes para as estruturas do poder presentes na sociedade \textbf{contemporânea}. \textbf{Enquanto} isso, a terceira competência trabalha as relações entre sociedade e natureza numa perspectiva socioambiental e de sustentabilidade, desenvolve o protagonismo do estudante, a reflexão, a ética socioambiental, o consumo responsável e a sustentabilidade global.
Já na quarta competência, o estudante ao analisar as relações de produção, capital e trabalho em diferentes territórios, contextos e culturas, discute o papel dessas relações na construção, consolidação e transformação das sociedades estimulando o seu protagonismo, e estabelece premissas a serem desenvolvidas em seu projeto de vida.
Ao se observar a quinta competência, podemos perceber \textbf{que} ela \textbf{opera} com visão e concepção de mundo, valores e atitudes \textbf{que} visam o combate às injustiças, \textbf{um} \textbf{compromisso} e respeito com a democracia e os Direitos Humanos.
A partir dela, é possível desenvolver o protagonismo dos estudantes na proposição de ações éticas de combate à violência e de formação ética.
E por fim, não menos importante, temos a sexta competência, \textbf{que} busca desenvolver a atuação protagonista do estudante mediante a participação pública e coletiva, com respeito à diversidade, fortalecendo a cidadania, promovendo o seu projeto de vida e a compreensão das demandas de povos indígenas e afrodescendentes no Brasil \textbf{atual}.
A BNCC das CHSA tem concepção marcada pela diversidade de juventudes e de narrativas, e, cada componente tem sua especificidade, mas há \textbf{um} diálogo intenso entre eles. \textbf{Vemos} a Filosofia no seu tratar do Ser Humano e dos Paradigmas do Pensamento, \textbf{enquanto} a Sociologia trata das Instituições e dos Grupos Sociais, a História dos processos e dinâmicas em diferentes Tempos e Espaços, e em Geografia as relações em Sociedade e Natureza, sempre em perspectiva interdisciplinar. \textbf{Uma} das características \textbf{fortes} da BNCC é o reforço a \textbf{uma} visão de área.
Porém com isso não quer \textbf{dizer} \textbf{que} os componentes curriculares estão dispensados.
Devemos sempre aferir as especificidades de cada componente devido às contribuições \textbf{que} cada \textbf{um} tem para os estudantes dessa etapa de ensino.
Neste sentido, devemos pensar na organização das \textbf{categorias} da área e o \textbf{que} cada \textbf{uma} delas possibilita na mediação do conhecimento escolar, \textbf{que} permite a melhor compreensão de mundo em \textbf{que} o jovem está inserido, realizando conexão com o seu projeto de vida.
Deste modo, as \textbf{categorias} das CHSA podem ser consideradas : > [... ] fundantes para a investigação e a aprendizagem, não se confundindo com temas ou propostas de conteúdo.
São aquelas cuja tradição nos diferentes campos das Ciências Humanas utiliza para a compreensão das ideias, dos fenômenos e dos processos políticos, sociais, econômicos e culturais.
Se, no Ensino Fundamental, essas \textbf{categorias} estão presentes na operacionalização das competências, habilidades e dos objetos de conhecimento, no Ensino Médio elas são \textbf{explicitadas} considerando a capacidade de abstração e simbolização dos estudantes.
Por sua vez, as competências e habilidades propostas permitem ampliar e aprofundar os conhecimentos já \textbf{sistematizados}, compreendendo -os em circunstâncias ( BRASIL, 2018, p. 550 ).
A primeira \textbf{categoria} é Tempo e Espaço, \textbf{que} esclarece os fenômenos nas CHSA ao identificar contextos, sendo \textbf{categorias} \textbf{que} sempre estão juntas e não se dissociam.
Na Etapa Ensino Médio, os objetos de conhecimento propiciam a análise de acontecimentos ocorridos em diferentes circunstâncias podendo fazer associações, comparações, assim como a compreensão dos processos marcados por continuidade, mudanças e rupturas.
Explicar o \textbf{que} é semelhante ou diferente em cada cultura não é \textbf{um} movimento \textbf{simples}.
É tão complexo quanto explicar as razões e os motivos ( materiais e imateriais ) responsáveis pela formação de \textbf{uma} sociedade, sua língua, usos e costumes.
Explicar a “ lógica ” \textbf{que} produz a diversidade humana seria considerado atividade difícil ( BRASIL, 2018, p. 563 ).
Nesse sentido, o próprio conceito de tempo e espaço pode adquirir diferentes significados, sendo considerado \textbf{um} desafio sobre o qual se debruçaram diferentes áreas do saber, como a Filosofia, a Matemática, a Biologia, e a própria área de CHSA.
O Tempo pode assumir significados diferentes e não é homogêneo, nem tampouco linear, salutar é a compreensão \textbf{que} o tempo assume significados e importância variados.
É salutar \textbf{que} na Etapa Ensino Médio os estudantes possam adquirir conhecimentos necessários \textbf{que} os levem à compreensão das diferentes dimensões da cronologia do tempo, sejam estas simbólicas ou abstratas, da concepção e noções do tempo e espaço para diferentes povos e sociedades.
O espaço não deve ser representado somente pelas cartografias, mas, sim, deve ser contemplado pelas dimensões históricas e culturais, estar associado aos diferentes arranjos de diversas naturezas \textbf{que} contemplem o ir e vir de povos e sociedades com seus eventos, disputas, tradições e ou denominações.
Assume também o espaço \textbf{um} lugar na distribuição e circulação das mercadorias dos insumos e consumo, realiza -se fluxos e estabelece -se relações de trocas e trabalho com seus ritmos variados.
A segunda \textbf{categoria} é Território e \textbf{Fronteira}, a qual pode ser considerada \textbf{bastante} ampla.
Na BNCC, essas duas \textbf{categorias} estão descritas da seguinte forma : > (... ) território é \textbf{uma} \textbf{categoria} usualmente associada a \textbf{uma} porção da superfície terrestre sob domínio de \textbf{um} grupo e suporte para nações, estados, países.
É dele \textbf{que} provêm alimento, segurança, identidade e refúgio.
Engloba as noções de lugar, região, \textbf{fronteira} e, especialmente, os limites políticos e administrativos de cidades, estados e países, sendo, portanto, esquemas abstratos de organização da realidade.
Associa -se território também à ideia de poder, jurisdição, administração e soberania, dimensões \textbf{que} expressam a diversidade das relações sociais e permitem juízos analíticos. \textbf{Fronteira} também é \textbf{uma} \textbf{categoria} construída \textbf{historicamente}.
Ao expressar \textbf{uma} cultura, povos definem \textbf{fronteiras}, formas de organização social e, por vezes, áreas de confronto com outros grupos.
A conformação dos impérios coloniais, a formação dos Estados Nacionais e os processos de globalização problematizam a discussão sobre limites culturais e \textbf{fronteiras} nacionais.
Os limites, por exemplo, entre civilização e barbárie geraram, não raro, a destruição daqueles indivíduos considerados bárbaros.
Temos aí \textbf{uma} \textbf{fronteira} sangrenta.
Povos com culturas e saberes distintos em muitos casos foram separados ou reagrupados de forma a resolver ou agravar conflitos, facilitar ou dificultar deslocamentos humanos, favorecer ou impedir a integração territorial de populações com identidades semelhantes ( BRASIL, 2018, p. 564 ).
Pode -se entender \textbf{que} os movimentos dentro das cidades são intencionais e denotam as marcações tradicionais dentro de \textbf{um} território, com suas formas de organizações sociais, \textbf{que} ampliam ou separam grupos dentro de suas especificidades e denotam saberes, cultura, arte, festas, musicalidades, tradições e lazer.
Desta forma, \textbf{demarcam} -se \textbf{fronteiras} \textbf{que} podem ser visíveis ou invisíveis e esferas do poder pertencentes a \textbf{um} território específico.
Quando se pensa nisso, pode -se explicar os diferentes saberes dentro de \textbf{fronteiras} \textbf{que} envolvem práticas de diferentes sociedades, \textbf{enquanto} habilidades como a caça, a pesca e ou a agricultura envolvem conhecimentos específicos e, em alguns momentos, ficam restritos a determinadas comunidades pelas próprias especificidades do cotidiano e podem estabelecer as dinâmicas pessoais ou mesmo coletivas de determinados indivíduos ou grupos, inseridos em espaços \textbf{que} os permitam desenvolver as competências dessas habilidades com seus conhecimentos e práticas.
Nesse sentido, a Etapa Ensino Médio é \textbf{uma} fase importante para \textbf{que} os estudantes se apropriem dessas \textbf{categorias} para a compreensão dos processos identitários marcados pelos diferentes Territórios e Fronteiras, para entender \textbf{que} a posição \textbf{que} ocupam no mundo é marcada por diferenciadas disputas e \textbf{que} cada \textbf{um} pode ajudar a transformar e mobilizar o lugar em \textbf{que} \textbf{vive}.
A terceira \textbf{categoria} agrupada em : Indivíduo, Natureza, Sociedade, Cultura e Ética, oferece ampla discussão de conceitos e noções, bem como, proporciona a \textbf{abordagem} de questões de tradição socrática, como por exemplo : ‘ O \textbf{que} é o ser humano? > Na busca da unidade, de \textbf{uma} natureza ( physis ), os primeiros pensadores gregos \textbf{sistematizaram} questões e se indagaram sobre as finalidades da existência, sobre o \textbf{que} era comum a todos os seres da mesma espécie, produzindo \textbf{uma} visão essencializada e metafísica sobre os seres humanos.
A identificação da condição humana como animal político – e animal social – significa \textbf{que}, independentemente da \textbf{singularidade} de cada \textbf{um}, as pessoas são essencialmente capazes de se organizar para \textbf{uma} vida em comum e de se governar.
Ou seja, os seres humanos têm \textbf{uma} necessidade vital da convivência coletiva. ( BRASIL, 2018, p.565 ).
Mas, para além de sua condição social e política, o ser humano tem a condição de subsistir, nesse viés, a busca constante pela sobrevivência.
Para isso, coexiste em \textbf{uma} ligação direta com a natureza no sentido de transformá -la, o \textbf{que} permite a interação com outros indivíduos, no agrupamento com outros, delineiam sua forma de vida, hábitos, costumes, tradição e valores como ser social, dentro de \textbf{um} grupo específico, apropriando -se e interagindo em \textbf{uma} determinada cultura.
Nas relações com demais indivíduos, passa a estabelecer ligações e interações sociais, constrói visão e compreensão do mundo, com signos, símbolos, significados e representações.
As diferentes formas de organização do espaço físico territorial presentes no mundo, com suas organizações próprias, culturas e tradições, apontam as formas diversas como esses grupos e povos se relacionam com a natureza, no sentido de protegê -la e preservar, bem como, dos problemas ambientais decorrentes do mau uso e mal cuidado da mesma, demonstra o grau cognitivo no \textbf{que} tange a sustentabilidade e a preservação da mesma.
Desse modo, as “ Ciências Humanas compreendem a cultura a partir de contribuições de diferentes campos do saber.
O caráter polissêmico da cultura permite compreender o modo como ela se apresenta a partir de códigos de comunicação e comportamento, de símbolos e artefatos ”, \textbf{que} representam parte de toda produção, circulação e consumo de sistemas culturais \textbf{que} se manifestam na vida social. ( BRASIL, 2018, p. 566 ).
Essas inserções dos indivíduos se apresentam em culturas mais diferenciadas possíveis, \textbf{que} vão desde a urbana, passando pela erudita, locais, regionais, entre outras.
Produzem conhecimentos e saberes \textbf{que} são geradoras de transformações nos campos culturais, sociais, políticos e econômicos do seu tempo.
Os novos códigos culturais presentes na sociedade, as concepções e as formas de organizações foram alterados na modernidade, devido à noção do indivíduo ter se tornado complexa pelas mudanças ocorridas, nas relações entre as pessoas, trazendo alterações no mundo ocidental na sua organização política.
Com isso, propiciou condições de repensar a natureza dos seres humanos com suas \textbf{singularidades}, capacidades e diversidades, mas sem esquecer os sistemas econômicos produtores de homogeneização de \textbf{sujeitos}, mas, nas relações sociais, produzem desigualdades.
Inferindo as diferenças e semelhanças entre os indivíduos em espaços e circunstâncias diversas, como no ambiente de trabalho, familiar, bairro, religião e outros, como também em relação a outros povos.
Importante garantir aos estudantes de Ciências Humanas e Sociais Aplicadas o olhar cuidadoso a essas diferenças, permitindo o entendimento de \textbf{que} os povos são diversos e devem coexistir em suas diversidades ao reconhecer os pontos de vistas e de mundo diferentes.
Desse modo, o aprofundamento das questões sociais, culturais e individuais permite enraizar, no Ensino Médio, a discussão sobre a ética.
A fim de permitir aos estudantes aprofundamentos sobre questões básicas \textbf{que} envolvem o respeito às diferenças, a convivência e o bem comum na perspectiva de pensar a ética dentre os princípios dos direitos humanos, como \textbf{estímulo} do acolhimento das diferenças entre diferentes povos com intuito de promover o convívio social e o respeito às pessoas, à coletividade e ao bem público.
Para tanto, tendo em vista o conhecimento do Outro, da outra cultura, isso possibilita o desenvolvimento da empatia ao se indagar e indagar ao Outro, atitude fundamental a ser “ desenvolvida na área de Ciências Humanas e Sociais Aplicadas.
Esse é o primeiro passo para a formação de \textbf{sujeitos} protagonistas tanto no processo de construção do conhecimento como da ação ética diante do mundo real e virtual, marcado por \textbf{uma} multiplicidade de culturas ”. ( BRASIL, 2018, p. 567 ).
A \textbf{categoria} Política e Trabalho também ocupa posição de destaque nas Ciências Humanas, \textbf{uma} vez \textbf{que} a vida em sociedade depreende práticas individuais e coletivas \textbf{que} são mediadas pela política e pelo trabalho.
Alguns dos temas \textbf{que} incitam a produção de conhecimento nessa área são : o diálogo sobre o bem comum e do público, dos regimes políticos e dos modelos de organização em sociedade, as lógicas de poder deliberadas em diferentes grupos, a micropolítica, as conjecturas a respeito do Estado e seus artifícios de legitimação, bem como a interferência da tecnologia nas formas de \textbf{sistematização} da sociedade.
A política está na gênese do pensamento filosófico na Grécia Antiga, quando a prática da argumentação e o diálogo a respeito dos destinos das cidades e suas leis impulsionaram a retórica ( in.
Rheloric. fr.
Rhétorique. ai.
Rhetorik it.
Retórica, ABBAGNANO, 2007, p. 856 ).
Arte de persuadir com o uso de instrumentos linguísticos e a abstração como técnicas indispensáveis para a discussão sobre o bem comum.
Nesse sentido, a política é compreendida como ação e inclusão do indivíduo na sociedade e no mundo, inserindo o \textbf{viver} coletivo e a cidadania.
Essa prática possibilitou ao cidadão da pólis perceber a política como \textbf{uma} edificação humana \textbf{que} é capaz de oportunizar as relações entre indivíduos e povos, bem como \textbf{fomentar} a crítica a procedimentos políticos, como a demagogia e a manipulação do interesse público.
Assim, a política é \textbf{uma} ferramenta empregada para lutar contra os autoritarismos, as tiranias, os terrores, as violências e as \textbf{inúmeras} maneiras de demolição da vida pública.
Na atualidade, as questões ligadas à política ganham maior visibilidade na geopolítica tanto em escala local quanto global, \textbf{uma} vez \textbf{que} apontam os conflitos entre pessoas, grupos, países e blocos transnacionais.
Desta forma, a política se \textbf{configura} como \textbf{uma} \textbf{categoria} \textbf{imprescindível} a ser examinada, \textbf{conhecida} e compreendida pelos estudantes.
Destarte, as discussões no \textbf{que} tange os modelos de \textbf{sistematização} do Estado, de governo e do poder são temáticas proferidas no Ensino Fundamental e aprofundadas e consolidadas no Ensino Médio, principalmente em seu aspecto formal e como sistemas jurídicos complexos.
Desta forma, são apresentados alguns elementos capazes de conciliar vários temas de ordem econômica, social, política, cultural e ambiental, \textbf{que} possibilitam principalmente o debate a respeito dos conceitos difundidos por diferentes sociedades e culturas.
A \textbf{categoria} trabalho, no \textbf{que} lhe concerne, abarca múltiplos aspectos – \textbf{antropológico}, filosófico, econômico, sociológico ou histórico : como virtude ; maneira de gerar riqueza, de controlar e de transformar a natureza ; como mercadoria ; ou como mecanismo de alienação.
Também é possível abordar a esfera trabalho como \textbf{categoria} examinada por diversos \textbf{autores}, como : trabalho como valor ( Karl Marx ) ; racionalidade capitalista ( Max Weber ) ; ou elemento de interação do indivíduo na sociedade em suas dimensões tanto corporativa como de integração social ( Émile Durkheim ).
O percurso ou os percursos definidos para desenvolver o tema deve \textbf{salientar} a \textbf{correlação} sujeito/trabalho, do mesmo modo \textbf{que} toda a sua teia de relações sociais.
Na atualidade, as mudanças na sociedade são amplas, principalmente por conta do emprego de novas tecnologias.
Desta forma, pode se perceber alterações nas maneiras de atuação dos trabalhadores nas diferentes esferas da produção, a diversificação das relações de trabalho, a variação índices de emprego e desemprego, o uso do trabalho inconstante, ou seja, temporário, a dispersão dos locais de trabalho, e o crescimento global da riqueza, suas diversas maneiras de concentração e \textbf{divisão}, e suas implicações a respeito das desigualdades sociais.
Existe \textbf{hoje} mais oportunidade para o empreendedorismo em todas as classes sociais, o \textbf{que} aumenta, desta forma, a relevância da educação financeira e do entendimento do sistema monetário nacional e mundial, \textbf{que} são primordiais para inclusão crítica e consciente no mundo \textbf{atual}.
Instituem -se, nessa conjuntura, novos desafios às Ciências Humanas, abrangendo o entendimento dos efeitos das inovações tecnológicas nas relações de produção, trabalho e consumo.
Assim, o conhecimento das \textbf{categorias} Política e Trabalho, ao longo do Ensino Médio, deve proporcionar aos estudantes examinar e entender a realidade, identificando os diversos projetos políticos e econômicos em disputa nas diferentes sociedades.
Por \textbf{conseguinte}, a heterogeneidade de concepções de mundo e a convivência com a diversidade propiciam ao estudante o desenvolvimento da sensibilidade, da autocrítica e da criatividade, resultando em ganhos éticos relacionados, ao protagonismo juvenil, isto é, à autonomia das deliberações em consonância com valores como liberdade, justiça social, pluralidade, solidariedade e sustentabilidade.
Isso \textbf{posto}, para desenvolver de forma significativa as \textbf{categorias} \textbf{inerentes} à área de Ciências Humanas, faz -se necessário o conhecimento da Base Nacional Comum Curricular e do Documento de Referência Curricular para Mato Grosso, etapas do Ensino Fundamental e Ensino Médio, para \textbf{que} as escolas possam reestruturar seu Projeto Político Pedagógico junto com a comunidade escolar.
Esse processo é de \textbf{suma} importância para a democracia e autonomia da escola, \textbf{uma} vez \textbf{que} o projeto pedagógico deve considerar as especificidades das escolas urbanas, do campo, quilombolas, indígenas e especializadas, possibilitando ao estudante : refletir criticamente sobre a realidade em \textbf{que} está inserido ; ser protagonista e solidário ; orientar a construção de seu projeto de vida e de sociedade ; usar os conhecimentos científicos, tecnológicos e \textbf{sócio-históricos} para resolver problemas e intervir criticamente nas questões sociais. [ ^24 ] : Em suas obras Edgar Morin evidencia a importância da composição de saberes, do constituir -se em relação ao próprio contexto social e o do estar no mundo percebendo \textbf{que} o todo é maior \textbf{que} a soma das partes.

% matched lemmas: abordagem, antropológico, aqui, atravessar, atual, autor, basear, bastante, categoria, compromisso, configurar, conhecido, conseguinte, contemporâneo, contextualização, correlação, demarcar, disciplina, divisão, dizer, embasar, enquanto, estímulo, explicitar, fomentar, forte, fronteira, historicamente, hoje, imprescindível, inerente, interdisciplinaridade, inúmero, método, novidade, operar, postar, que, salientar, simples, singularidade, sistematizar, sistematização, sujeito, sumo, sócio-histórico, territorialidade, um, ver, viver
\end{textsample}
